\documentclass[10pt]{article}
\usepackage[latin9]{inputenc}
\usepackage[a4paper]{geometry}
\geometry{verbose,tmargin=2cm,bmargin=2cm,lmargin=2cm,rmargin=2cm}
\usepackage{xcolor}
\usepackage{float}
\usepackage{calc}
\usepackage{bm}
\usepackage{amsmath}
\usepackage{amssymb}
\usepackage{graphicx}
\PassOptionsToPackage{normalem}{ulem}
\usepackage{ulem}

\makeatletter
%%%%%%%%%%%%%%%%%%%%%%%%%%%%%% User specified LaTeX commands.
%!TEX program = xelatex
\usepackage{mathrsfs}
\usepackage{titlesec}
\usepackage{xcolor}
\usepackage{enumerate}
\usepackage{bm}
\usepackage{tikz}
\usepackage{listings}
\usetikzlibrary{arrows}
\usepackage{subfigure}
\usepackage{graphicx}
\usepackage{booktabs}\usepackage{multirow}\usepackage{upquote}
\usepackage{float}
\usepackage{pdfpages}




\lstset{language=Matlab}
\lstset{breaklines}

\input defs.tex
\newtheorem{definition}{Definition}
\newtheorem{proposition}{Proposition}\newtheorem{remark}{Remark}
\newcommand{\vectdot}[1]{\dot{\mathbf{#1}}}
\titleformat*{\section}{\centering\LARGE\scshape}
\renewcommand{\thesection}{\Roman{section}}
\lstset{language=Matlab,tabsize=4,frame=shadowbox,basicstyle=\footnotesize,
keywordstyle=\color{blue!90}\bfseries,breaklines=true,commentstyle=\color[RGB]{50,50,50},stringstyle=\ttfamily,numbers=left,numberstyle=\tiny,
  numberstyle={\color[RGB]{192,92,92}\tiny},backgroundcolor=\color[RGB]{245,245,244},inputpath=code}

\makeatother

\begin{document}
\title{CS182: Introduction to Machine Learning, Fall 2025\\
 Homework 1}
\author{{\small (Due on Sunday, Oct. 19 at 11:59pm (CST))}}
\date{\date{}}

\maketitle
{\small{}%
\noindent\fcolorbox{black}{white}{\begin{minipage}[t]{1\columnwidth - 2\fboxsep - 2\fboxrule}%
Notice:
\begin{itemize}
\item Please submit your assignments via Gradescope. The entry code is \uline{428YBY}.
\item Please make sure you select your answer to the corresponding question
when submitting your assignments.
\item Each person has a total of five days to be late without penalty for
all the assignments. Each late delivery less than one day will be
counted as one day.
\end{itemize}
%
\end{minipage}}}{\small\par}
\begin{enumerate}
\item \defpoints{20} {[}\textit{Linear Algebra}{]} In robotics, the Jacobian Matrix, $\mathbf{J}$, is a fundamental tool that creates a linear mapping between the robot's joint velocities, $\vectdot{q}$, and its end-effector velocity, $\mathbf{v}$, through the relationship: $\mathbf{v} = \mathbf{J}\vectdot{q}$. By analyzing the properties of the Jacobian, we can evaluate a robot's dexterity and motion capability in a given configuration, an analysis known as manipulability.


\begin{itemize}
\item[(a)]  For a robot in a specific configuration, its Jacobian is given by the matrix $\mathbf{J} = \begin{pmatrix} 1 & 0 \\ 0 & 1 \\ 1 & 1 \end{pmatrix}$. In manipulability analysis, the symmetric matrix $\mathbf{S} = \mathbf{J}^T \mathbf{J}$ is of critical importance, as its eigenvalues and eigenvectors reveal the intrinsic relationship between joint motions and end-effector motions.

\begin{enumerate}
    \item Calculate the symmetric matrix $\mathbf{S} = \mathbf{J}^T \mathbf{J}$ for this robot.\defpoints{3}
    \item Find all eigenvalues and their corresponding orthonormal eigenvectors for the matrix $\mathbf{S}$.\defpoints{3}
\end{enumerate}

\item[(b)]  The Singular Value Decomposition (SVD) of the Jacobian is given by $\mathbf{J} = \mathbf{U}\boldsymbol{\Sigma} \mathbf{V}^T$. The singular values, $\sigma_i$, on the diagonal of $\boldsymbol{\Sigma}$ are a direct measure of the robot's ability to move in different directions.
Prove that the squared singular values of the Jacobian $\mathbf{J}$ are the eigenvalues of the manipulability matrix $\mathbf{J}^T \mathbf{J}$. \defpoints{6}

\item[(c)] The squared magnitude of the end-effector velocity is given by $\|\mathbf{v}\|^2 = \mathbf{v}^T \mathbf{v} = (\mathbf{J}\vectdot{q})^T(\mathbf{J}\vectdot{q}) = \vectdot{q}^T \mathbf{J}^T \mathbf{J} \vectdot{q}$. To plan the most efficient movements, engineers need to determine the maximum possible end-effector speed when the joints move with a unit norm speed (i.e., $\|\vectdot{q}\|^2 = \vectdot{q}^T\vectdot{q} = 1$).
Prove that for any unit-norm joint velocity vector $\vectdot{q}$, the maximum value of the squared end-effector velocity, $\|\mathbf{v}\|^2$, is equal to the largest eigenvalue of the matrix $\mathbf{J}^T \mathbf{J}$.\defpoints{8}
\end{itemize}

\newpage

% \begin{enumerate}
\item \defpoints{20} {[}\textit{Linear Algebra}{]} 
Let $f(\mathbf{x}) = \frac{1}{2}\mathbf{x}^T\mathbf{A}\mathbf{x} - \boldsymbol{b}^T\mathbf{x} + c$, where $\mathbf{A} \in \mathbb{R}^{n\times n}$ is symmetric, $\boldsymbol{b} \in \mathbb{R}^n$, and $c \in \mathbb{R}$.
\begin{itemize}
	\item[(a)] Derive the Gradient and Hessian of $f(\mathbf{x})$. \defpoints{8}
    
	\item[(b)] Under what condition on $\mathbf{A}$ does $f(\mathbf{x})$ correspond to a minimum, and find the point $\mathbf{x}^\ast$ that makes $f(\mathbf{x})$ smallest. \defpoints{6} %Show that the critical point $\mathbf{x}^\ast = \mathbf{A}^{-1}\boldsymbol{b}$ minimizes $f(\mathbf{x})$ if and only if $A$ is positive definite. \defpoints{6}
    
	\item[(c)] Suppose $\mathbf{A}$ has eigenvalues $\lambda_1, \lambda_2$. Describe the general shape of $f(\mathbf{x})=\frac{1}{2} \mathbf{x}^T\mathbf{A}\mathbf{x}$ in 2D for the following cases. 
    \begin{itemize}
        \item[(1)] $\lambda_1, \lambda_2 \textgreater 0.$ \defpoints{2}
        \item[(2)] $\lambda_1, \lambda_2 \textless 0.$ \defpoints{2}
        \item[(3)] $\lambda_1\textgreater 0, \lambda_2 \textless 0.$ \defpoints{2}
    \end{itemize} %Using the eigendecomposition $A = Q\Lambda Q^T$, interpret geometrically how the curvature (Hessian) scales along each eigenvector direction. \defpoints{6}
\end{itemize}


\newpage{}




\item\defpoints{20} {[}\textit{Probability and Statistics}{]}
\begin{itemize}
\item[(a)] For any non-negative random variable $X$ with finite variance, prove that
     \[\Pr[X>0]\geq \frac{(\textrm{E}[X])^2}{\textrm{E}[X^2]}.\]
Hint: Just use the fact that the expectation \(\mathrm{E}[UV]\) is a valid inner product of random variables \(U\) and \(V\) and you can therefore use the Cauchy-Schwarz inequality.
\defpoints{4}
\item[(b)] For a random variable X with finite variance and E[X] = 0, prove that
     \[\Pr[X\geq t]\leq \frac{\textrm{E}[X^2]}{\textrm{E}[X^2]+t^2}\text{ for any }t\geq 0\]
     Hint: You can try the same logic as in the previous question.\defpoints{6}
\item[(c)] Prove that the following definition of Gaussian random vector is equivalent to the definition via the density function:
\begin{definition}\label{gaussian-as-affine-map}
    Gaussian random vectors are affine transformations of IID gaussian random variables. In other words, if X has nonsingular $\mathbf{\Sigma_X}$, then we can write
\[\mathbf{X = \mu_x +AW}\]for some full rank \(\mathbf A\in\mathbb{R}^{n\times n}\) and \(\mathbf W = (W_1,..., W_n)^T\), where \(W_i\sim_{\mathrm{IID}}\mathcal{N} (0, 1)\) 
\end{definition}
Hint: For strictly monotone and differentiable function $g(x)$ such that $Y=g(X)$, and let \(g^{-1}\) be the inverse function of \(g\)\[f_Y(y)=f_X(g^{-1}(y))\frac{1}{|g'(g^{-1}(y)|}\]
\defpoints{10}
\end{itemize}
\newpage{}



\item \defpoints{20} {[}\textit{Bayesian Decision Theory}{]} A major technology
company is hiring for a Machine Learning Engineer position and uses
a two-stage screening process: 
\begin{itemize}
\item 1. Resume screening: (binary outcome: $\textit{Pass}$ or $\textit{Fail}$)
\item 2. Technical interview (binary outcome: $\textit{Hire}$ or $\textit{Reject}$)
\end{itemize}
Historical hiring data reveals the following:
\begin{itemize}
\item 20\% of all applicants are ultimately qualified for the role. (i.e.,
$P\left(\omega_{Q}\right)=0.2,$ where $\omega_{Q}=\mathrm{qualified}$,
$\omega_{U}=\mathrm{unqualified}$).
\item Among qualified candidates, 90\% pass the resume screen: $P\left(\mathrm{Pass}\mid\omega_{Q}\right)=0.9$.
\item Among unqualified candidates, 30\% pass the resume screen: $P\left(\mathrm{Pass}\mid\omega_{U}\right)=0.3$.
\item The company incurs different costs for hiring decisions:
\begin{itemize}
\item Correctly hiring a qualified candidate yields a benefit of +10 (i.e.,
cost = --10).
\item Rejecting a qualified candidate has a cost of 5 (missed opportunity)
\item Hiring an unqualified candidate has a cost of 15 (training, turnover,
etc.).
\item Correctly rejecting an unqualified candidate has zero cost.
\end{itemize}
\item[(a)]  Compute the posterior probability that an applicant who passed the
resume screen is actually qualified, i.e. $P\left(\omega_{Q}\mid\mathrm{Pass}\right)$.\defpoints{4}
\item[(b)]  Using Bayesian decision theory, determine whether the company should
automatically hire every applicant who passes the resume screen, or
reject them, based on minimizing expected risk. Justify your answer
by computing the conditional risk for both actions.\defpoints{4}
\item[(c)]  Suppose the company introduces a new AI-powered resume screener
that improves the true positive rate to 95\% while reducing the false
positive rate to 10\%. Recalculate the posterior $P\left(\omega_{Q}\mid\mathrm{Pass}\right)$
and the optimal decision in question (b).
\defpoints{8}
\item[(d)] Discuss how advancements in screening technology affect hiring decisions
from a Bayesian perspective.\defpoints{4}
\end{itemize}

\newpage


\item \defpoints{20} {[}\textit{Bayesian Decision Theory}{]}
		Given a training dataset $T = \{(\mathbf{x}_1, y_1), (\mathbf{x}_2, y_2), \ldots, (\mathbf{x}_N, y_N)\}$, where $\mathbf{x}_i = (x_i^{(1)}, x_i^{(2)}, \ldots, x_i^{(n)})$ is a feature vector and $y_i \in \{c_1, c_2, \ldots, c_K\}$ is a class label. Bayesian estimation often add a small constant $\lambda$ to each count, i.e., 
		\[
		P_{\lambda}(Y = c_k) = \frac{\sum_{i=1}^N I(y_i = c_k) + \lambda}{N + K \lambda},\quad P_{\lambda}(X^{(j)} = a_{jl} | Y = c_k) = \frac{\sum_{i=1}^N I(x_i^{(j)} = a_{jl}, y_i = c_k) + \lambda}{\sum_{i=1}^N I(y_i = c_k) + S_j \lambda}.
		\]
		where $S_j$ is the number of possible values for feature $j$.
		
		\begin{table}[h]
			\centering
			\caption{Dataset for playing game}
			\label{tab:dataset}
			\begin{tabular}{ccc}
				\toprule
				Outlook & Temperature & Play \\
				\midrule
				Sunny & Hot & No \\
				Sunny & Mild & No \\
				Rainy & Hot & Yes \\
				Rainy & Mild & Yes \\
				\bottomrule
			\end{tabular}
		\end{table}
		
		Consider the dataset \( T \) in Table~\ref{tab:dataset} which records whether a game is played based on the outlook and temperature. The features are: Outlook: Sunny (S) or Rainy (R), Temperature (Temp): Hot (H) or Mild (M). The class label Play include: Yes (Y) or No (N). 
		
		\begin{itemize}
			\item[(a)] Estimate the prior probabilities \( P(\text{Play}) \) and the conditional probabilities \( P(\text{Outlook} | \text{Play}) \) and \( P(\text{Temp} | \text{Play}) \) using maximum likelihood estimation. Then, for a new instance with \text{Outlook}=\text{R}ainy and \text{Temp}=\text{H}, compute the posterior probabilities \( P(\text{Play}=\text{Y} | \text{Outlook}=\text{R}, \text{Temp}=\text{H}) \) and \( P(\text{Play}=\text{N} | \text{Outlook}=\text{R}, \text{Temp}=\text{H}) \) using the naive Bayes classifier. \defpoints{10}
			\item[(b)] Use Bayesian estimation with Laplacian smoothing (with smoothing parameter \( \lambda = 1 \)). Recompute the prior probabilities (with smoothing) and the conditional probabilities. Then, compute the posterior probabilities for the same new instance. \defpoints{5}
			\item[(c)] Based on the results of the above calculations, discuss why Laplacian smoothing is used. \defpoints{5}
		\end{itemize}
\end{enumerate}





\end{document}
